% GENERATED_BY: run_oc_core_monograph_translation_rework_dispatch_v1.py
% AI_ASSISTED_TRANSLATION_DRAFT: TRUE
% THEOREM_REVIEW_STATUS: APPROVED
% THEOREM_REVIEW_MODE: FORMAL_AI_GATE__CODEX_CLI_CHATGPT
% THEOREM_REVIEW_REVIEWER_ID: CODEX_CLI_CHATGPT__AUTO_DISPATCH_20260406
% THEOREM_REVIEW_AUTHORITY_CLASS: FINAL_ADVERSARIAL_EXTERNAL_LLM_REVIEWER
% THEOREM_REVIEWED_AT: 2026-04-14T11:19:57Z
% SOURCE_LANGUAGE: EN
% TARGET_LANGUAGE: DE
% SOURCE_EN_REF: appendix/E_oc_core_1_3_journal_core_bridge.tex
\section{Brücke zur Journal-Kern-Extraktion}

Dieser Anhang erläutert, wie der Journal-Kern-Artikel mit der Hauptmonographie zusammenhängt.

\subsection{Warum der Journal-Kern existiert}

Einige Kanäle erfordern ein abgegrenztes, schneller lesbares wissenschaftliches Artefakt.
Herausgeber, Gutachter und externe Leser benötigen oft ein Objekt im Umfang eines Artikels, dessen verteidigte Behauptung
in einem Durchgang gelesen werden kann.
Der Journal-Kern existiert zu diesem Zweck.
Seine Aufgabe ist nicht, die Monographie zu ersetzen, sondern einen eng gefassten Artikel bereitzustellen, der an
dieselben wissenschaftlichen Wahrheitsbedingungen gebunden bleibt.

\subsection{Was der Journal-Kern weglassen darf}

Der Journal-Kern darf verdichten:

\begin{itemize}
\item umfangreiche didaktische Erläuterungen auf Kapitelebene,
\item große Teile des Domänenatlas,
\item breite Modulabdeckung,
\item detaillierte Inventare in den Anhängen,
\item und einen Teil des pädagogischen Gerüsts auf Monographieebene.
\end{itemize}

Er darf nicht stillschweigend ändern:

\begin{itemize}
\item den Geltungsbereich des Theorems,
\item explizite Nichtbehauptungen,
\item die beanspruchte Quellenbasis,
\item die Bedeutung von Abbildungen,
\item oder die Beziehung zwischen dem geerbten Korpus und der gegenwärtig verteidigten Behauptung.
\end{itemize}

\subsection{Was die Monographie späteren Artefakten garantiert}

Weil der Journal-Kern aus dem vorliegenden Band extrahiert wird, müssen spätere eigentümergerichtete und öffentliche
Kanäle ihr eigenes Verständnis der Wissenschaft nicht improvisieren.
Sie können aus der Hauptmonographie ableiten, wenn sie Tiefe benötigen, aus dem Journal-Kern, wenn sie eine
abgegrenzte Darstellung benötigen, und aus der Begleitebene, wenn sie stützende Belege benötigen.

Das ist der praktische Publikationswert der Monographie.
Sie ist nicht bloß größer als der Artikel.
Sie ist das Dokument, das verhindert, dass jede spätere Projektion zu einer eigenständigen
Interpretationsübung wird.

\clearpage
